%%%%%%%%%%%%%%%%%%%%%%%%%%%%%%%%%%%%%%%%%%%%%%%%%%%%%%%%%%%%%%%%%%%%%%%%%%%%%%%%
%2345678901234567890123456789012345678901234567890123456789012345678901234567890
%        1         2         3         4         5         6         7         8

\documentclass[letterpaper, 10 pt, conference]{ieeeconf}  % Comment this line out if you need a4paper

%\documentclass[a4paper, 10pt, conference]{ieeeconf}      % Use this line for a4 paper

\IEEEoverridecommandlockouts                              % This command is only needed if 
                                                          % you want to use the \thanks command

\overrideIEEEmargins                                      % Needed to meet printer requirements.

%In case you encounter the following error:
%Error 1010 The PDF file may be corrupt (unable to open PDF file) OR
%Error 1000 An error occurred while parsing a contents stream. Unable to analyze the PDF file.
%This is a known problem with pdfLaTeX conversion filter. The file cannot be opened with acrobat reader
%Please use one of the alternatives below to circumvent this error by uncommenting one or the other
%\pdfobjcompresslevel=0
%\pdfminorversion=4

% See the \addtolength command later in the file to balance the column lengths
% on the last page of the document

% The following packages can be found on http:\\www.ctan.org
%\usepackage{graphics} % for pdf, bitmapped graphics files
%\usepackage{epsfig} % for postscript graphics files
%\usepackage{mathptmx} % assumes new font selection scheme installed
%\usepackage{times} % assumes new font selection scheme installed
%\usepackage{amsmath} % assumes amsmath package installed
%\usepackage{amssymb}  % assumes amsmath package installed

\usepackage[utf8]{inputenc}
\usepackage[T1]{fontenc}
\usepackage{amsmath,amssymb}
\usepackage{graphicx}
\usepackage{booktabs}
\usepackage{enumitem}
\usepackage{xcolor}
\usepackage{cite}
\usepackage{hyperref}
\usepackage{url}
\usepackage{microtype}
\usepackage{algorithm}
\usepackage{algorithmic}
\usepackage{multirow}
\usepackage{balance}
\usepackage{tabularx}
\usepackage{array}


\title{\LARGE \bf
Preparation of Papers for IEEE Sponsored Conferences \& Symposia*
}


\author{Albert Author$^{1}$ and Bernard D. Researcher$^{2}$% <-this % stops a space
\thanks{*This work was not supported by any organization}% <-this % stops a space
\thanks{$^{1}$Albert Author is with Faculty of Electrical Engineering, Mathematics and Computer Science,
        University of Twente, 7500 AE Enschede, The Netherlands
        {\tt\small albert.author@papercept.net}}%
\thanks{$^{2}$Bernard D. Researcheris with the Department of Electrical Engineering, Wright State University,
        Dayton, OH 45435, USA
        {\tt\small b.d.researcher@ieee.org}}%
}


\begin{document}



\maketitle
\thispagestyle{empty}
\pagestyle{empty}


%%%%%%%%%%%%%%%%%%%%%%%%%%%%%%%%%%%%%%%%%%%%%%%%%%%%%%%%%%%%%%%%%%%%%%%%%%%%%%%%
\begin{abstract}

This electronic document is a ÒliveÓ template. The various components of your paper [title, text, heads, etc.] are already defined on the style sheet, as illustrated by the portions given in this document.

\end{abstract}


%%%%%%%%%%%%%%%%%%%%%%%%%%%%%%%%%%%%%%%%%%%%%%%%%%%%%%%%%%%%%%%%%%%%%%%%%%%%%%%%
\section{INTRODUCTION}

Camera calibration estimates camera intrinsics and target poses from
image observations of a target with known geometry. Planar checkerboards
and AprilGrids are widely used because of their simple structures and
reliable feature detection. For conventional lenses, translating and
rotating a single planar target usually provides sufficient observations
to constrain the camera model. Calibration becomes more challenging for
wide-angle cameras, whose projection models are highly nonlinear and
more sensitive to initialization and observation conditions. Recent
analyses have shown that observations confined to a limited radial span
can couple focal scale with projection-shape parameters, resulting in
poorly separated parameter directions and ill-conditioned linearized
updates~\cite{cocalib}. Measurements spanning broader radial regions, in
contrast, improve their local separability. Wide-angle calibration
therefore requires not only a sufficient number of features, but also
reliable observations across different radial regions.

However, realizing broad radial coverage with a single planar target is
difficult in practice. In each view, a single target occupies only a
limited and continuous image region, and its radial coverage is bounded
by its projected extent. Enlarging the target projection increases the
radial span, but also forces the same pattern to cross image regions
with substantially different projection scales and distortion levels,
placing greater demands on feature detection and subpixel localization.
Existing wide-angle calibration methods mitigate this problem by using
an intermediate camera model to reproject the target and adaptively
refine missing or distorted features. Nevertheless, a practical trade-off
remains between radial coverage and local observation reliability. This
motivates a calibration framework that distributes reliable observations
across different field-of-view regions while preserving stable local
feature detection.

Motivated by these findings, we construct a distributed multi-board
target using multiple boards, each carrying a uniquely coded AprilTag.
The boards can be independently detected and associated within the same
image, allowing observations to be distributed across different spatial
and radial regions without requiring a single large continuous target to
span the strongly distorted image domain. Each tag directly provides
only four sparse observations, referred to as outer corners. Although
these distributed outer corners are sufficient to bootstrap an
intermediate camera model, they remain locally sparse for final
calibration. We therefore use the intermediate model to predict the
image locations of predefined internal lattice points and refine them
using local image evidence. In this way, spatially distributed but
locally sparse outer-corner observations are extended into dense and
geometrically consistent observations across different field-of-view
regions. Based on this design, we develop a multi-board calibration
system for wide-angle central cameras that integrates distributed target
detection, dense observation recovery, and camera calibration into a
unified pipeline.

Our main contributions are summarized as follows:

\begin{itemize}
\setlength{\itemsep}{1pt}
\setlength{\parskip}{0pt}
\setlength{\parsep}{0pt}
\item We present an open-source easy-to-use multi-board calibration system designed for
wide-angle cameras.
\item To address insufficient spatial coverage of calibration observations
and unstable feature refinement, we propose a distributed multi-board design based on single AprilTags, together with an intermediate-model-assisted internal lattice-point refinement method.
The independently detectable single-tag boards extend observation coverage,
while the recovered and refined internal lattice observations provide dense
calibration constraints across different regions of the field of view.
\item 
\end{itemize}

Tests on several datasets and camera models cover reprojection error, ablations,
cross-backend calibration, and held-out pose estimation. The results show
reliable calibration and competitive accuracy.

\section{}

\subsection{Selecting a Template (Heading 2)}

First, confirm that you have the correct template for your paper size. This template has been tailored for output on the US-letter paper size. 
It may be used for A4 paper size if the paper size setting is suitably modified.

\subsection{Maintaining the Integrity of the Specifications}

The template is used to format your paper and style the text. All margins, column widths, line spaces, and text fonts are prescribed; please do not alter them. You may note peculiarities. For example, the head margin in this template measures proportionately more than is customary. This measurement and others are deliberate, using specifications that anticipate your paper as one part of the entire proceedings, and not as an independent document. Please do not revise any of the current designations

\section{Method}

\subsection{Multi-Board Target Formulation}

Our calibration target consists of $N$ rigid boards, each carrying a
uniquely coded AprilTag:
\begin{equation}
\mathcal{B}=\{B_0,B_1,\ldots,B_{N-1}\},
\end{equation}
where $B_0$ is selected as the reference board. Let
$\mathcal{F}_{B_b}$ denote the local coordinate frame of board $B_b$,
and let $\mathbf{p}_{b,j}\in\mathbb{R}^3$ denote the $j$-th target point
expressed in $\mathcal{F}_{B_b}$. The target points include the directly
detected outer corners and the internal lattice points recovered in
Sec.~\ref{sec:internal_point_recovery}.

We use $\mathcal{F}_{B_0}$ as the common coordinate frame of the
multi-board target and fix
$\mathbf{T}_{B_0B_0}=\mathbf{I}$.
Let $\mathbf{T}_{B_0B_b}\in SE(3)$ denote the static rigid transform
from $\mathcal{F}_{B_b}$ to $\mathcal{F}_{B_0}$. For frame $t$, let
$\mathcal{F}_{C_t}$ be the camera coordinate frame and let
$\mathbf{T}_{C_tB_0}\in SE(3)$ denote the pose of the reference board
in $\mathcal{F}_{C_t}$. A control point on board $B_b$ is therefore
expressed in the camera frame as
\begin{equation}
\mathbf{x}_{t,b,j}
=
\mathbf{T}_{C_tB_0}
\mathbf{T}_{B_0B_b}
\mathbf{p}_{b,j}.
\end{equation}

Given the camera parameters $\boldsymbol{\theta}$ and the projection
function $\pi(\cdot;\boldsymbol{\theta})$, the predicted image
location is
\begin{equation}
\hat{\mathbf{u}}_{t,b,j}
=
\pi\!\left(
\mathbf{x}_{t,b,j};
\boldsymbol{\theta}
\right)
=
\pi\!\left(
\mathbf{T}_{C_tB_0}
\mathbf{T}_{B_0B_b}
\mathbf{p}_{b,j};
\boldsymbol{\theta}
\right).
\end{equation}
The corresponding image observation is denoted by
$\mathbf{u}_{t,b,j}\in\mathbb{R}^2$. We collect all valid
observations in
\begin{equation}
\mathcal{O}
=
\left\{
(t,b,j,\mathbf{u}_{t,b,j})
\right\}.
\end{equation}
The unknown variables comprise the camera parameters
$\boldsymbol{\theta}$, the per-frame reference-board poses
$\{\mathbf{T}_{C_tB_0}\}$, and the static inter-board transforms
$\{\mathbf{T}_{B_0B_b}\}_{b=1}^{N-1}$.

\subsection{Camera Model Initialization}

At the beginning of calibration, no reliable camera model is available,
while a usable initialization is important for nonlinear wide-angle
calibration. The distributed AprilTag boards provide robust outer-corner
observations across different image and radial regions. We use these
observations to initialize the camera model. For a given camera-model
family, we construct an initial parameter vector
$\boldsymbol{\theta}^{0}$ using the image resolution and model-specific
default values.

For board $B_b$ observed in frame $t$, let
$\mathcal{C}^{\mathrm{out}}_{t,b}$ denote its valid outer-corner
observations. We define the reprojection residual as
\begin{equation}
\mathbf{r}^{\mathrm{out}}_{t,b,j}
(\boldsymbol{\theta},\mathbf{T})
=
\mathbf{u}_{t,b,j}
-
\pi\!\left(
\mathbf{T}\mathbf{p}_{b,j};
\boldsymbol{\theta}
\right).
\end{equation}
Given $\boldsymbol{\theta}^{0}$, a temporary frame--board pose is
initialized by
\begin{equation}
\widetilde{\mathbf{T}}^{0}_{C_tB_b}
=
\operatorname*{arg\,min}_{\mathbf{T}\in SE(3)}
\sum_{j\in\mathcal{C}^{\mathrm{out}}_{t,b}}
\left\|
\mathbf{r}^{\mathrm{out}}_{t,b,j}
(\boldsymbol{\theta}^{0},\mathbf{T})
\right\|_2^2 .
\end{equation}

We perform this initialization independently for all valid frame--board
observations and reject estimates with invalid poses or excessive
reprojection errors. Let $\mathcal{S}$ denote the retained frame--board
pairs. We then jointly refine the camera parameters and temporary poses:
\begin{equation}
\begin{aligned}
&\boldsymbol{\theta}^{*},
\left\{
\widetilde{\mathbf{T}}^{*}_{C_tB_b}
\right\}
\\
&\quad =
\operatorname*{arg\,min}_{
\boldsymbol{\theta},
\{\widetilde{\mathbf{T}}_{C_tB_b}\}}
\sum_{(t,b)\in\mathcal{S}}
\sum_{j\in\mathcal{C}^{\mathrm{out}}_{t,b}}
\rho\!\left(
\left\|
\mathbf{r}^{\mathrm{out}}_{t,b,j}
\left(
\boldsymbol{\theta},
\widetilde{\mathbf{T}}_{C_tB_b}
\right)
\right\|_2^2
\right).
\end{aligned}
\end{equation}
where $\rho(\cdot)$ is a robust loss function. At this stage, the boards
are not coupled through the unknown inter-board geometry; each valid
frame--board observation is assigned an independent pose. The optimized
parameters $\boldsymbol{\theta}^{*}$ initialize the subsequent complete
multi-board calibration and serve as the intermediate camera model for
internal-point recovery.

\subsection{Internal Lattice Point Recovery}
\label{sec:internal_point_recovery}

The outer corners provide stable board localization and are sufficient
to bootstrap an intermediate camera model. However, the four
observations from each tag remain locally sparse for final calibration.
We therefore use the intermediate model to recover image observations
for predefined internal lattice points, combining model-guided
prediction with local image evidence.

Given the intermediate camera parameters $\boldsymbol{\theta}^{*}$,
we reuse the temporary frame--board pose
$\widetilde{\mathbf{T}}^{*}_{C_tB_b}$ obtained in Sec.~\ref{sec:camera_init}.
For an internal lattice point
$\mathbf{p}_{b,j}$, $j\in\mathcal{C}^{\mathrm{int}}_b$, its direct
geometric projection
\begin{equation}
\widehat{\mathbf{u}}^{\mathrm{nom}}_{t,b,j}
=
\pi\!\left(
\widetilde{\mathbf{T}}^{*}_{C_tB_b}
\mathbf{p}_{b,j};
\boldsymbol{\theta}^{*}
\right)
\label{eq:nominal_internal_prediction}
\end{equation}
provides a nominal prediction and is used as a fallback when the
boundary-based seed described below is unavailable.

To construct the primary seed, we sample support points along the four
observed tag boundaries and back-project them onto the unit sphere using
the intermediate camera model. The resulting top, bottom, left, and
right boundary-ray curves are denoted by
$\mathbf{r}_{T}(s)$, $\mathbf{r}_{B}(s)$,
$\mathbf{r}_{L}(t)$, and $\mathbf{r}_{R}(t)$, respectively.
Let $(s_j,t_j)\in[0,1]^2$ be the normalized horizontal and vertical
coordinates of $\mathbf{p}_{b,j}$ on the board lattice. We interpolate
the opposite boundary curves as
\begin{equation}
\begin{aligned}
\mathbf{r}_{v}
&=
\operatorname{slerp}\!\left(
\mathbf{r}_{T}(s_j),
\mathbf{r}_{B}(s_j),
t_j
\right),\\
\mathbf{r}_{h}
&=
\operatorname{slerp}\!\left(
\mathbf{r}_{L}(t_j),
\mathbf{r}_{R}(t_j),
s_j
\right).
\end{aligned}
\label{eq:spherical_boundary_interpolation}
\end{equation}
The two interpolated rays are blended and projected back to the image:
\begin{equation}
\widetilde{\mathbf{d}}^{\mathrm{int}}_{t,b,j}
=
\frac{\mathbf{r}_{v}+\mathbf{r}_{h}}
{\left\|\mathbf{r}_{v}+\mathbf{r}_{h}\right\|_2},
\qquad
\widetilde{\mathbf{u}}^{\mathrm{int}}_{t,b,j}
=
\pi\!\left(
\widetilde{\mathbf{d}}^{\mathrm{int}}_{t,b,j};
\boldsymbol{\theta}^{*}
\right).
\label{eq:internal_point_seed}
\end{equation}

Starting from $\widetilde{\mathbf{u}}^{\mathrm{int}}_{t,b,j}$, we
perform a local search in the viewing-ray domain and then refine the
candidate to subpixel accuracy in the original image. A recovered
observation is rejected if it falls outside the image, lacks sufficient
local image evidence, or deviates from its seed by more than a predefined
threshold. Each accepted observation retains its board and lattice-point
indices and is added to $\mathcal{O}$ for joint optimization together
with the outer corners.

\subsection{Spherical Bundle Adjustment}

Standard bundle adjustment minimizes reprojection errors in the image
plane. For wide-angle cameras, however, the ray deviation corresponding
to a pixel residual depends on the image location and the local Jacobian
of the projection model. We therefore provide optional ray-domain and
hybrid objectives for final refinement.

Given an image observation $\mathbf{u}_{t,b,j}$, we back-project it
using the current camera parameters and normalize it to obtain the
observed ray:
\begin{equation}
\mathbf{d}^{\mathrm{obs}}_{t,b,j}
=
\frac{
\pi^{-1}(\mathbf{u}_{t,b,j};
\boldsymbol{\theta})
}{
\left\|
\pi^{-1}(\mathbf{u}_{t,b,j};
\boldsymbol{\theta})
\right\|_2
}.
\end{equation}
The corresponding predicted ray is obtained from the transformed target
point defined in Sec.~IV-A:
\begin{equation}
\mathbf{d}^{\mathrm{pred}}_{t,b,j}
=
\frac{\mathbf{x}_{t,b,j}}
{\|\mathbf{x}_{t,b,j}\|_2}.
\end{equation}
Let $\mathbf{B}_{t,b,j}\in\mathbb{R}^{3\times2}$ be an orthonormal
tangent basis at $\mathbf{d}^{\mathrm{obs}}_{t,b,j}$. The local ray
residual is defined as
\begin{equation}
\mathbf{r}^{\mathrm{ray}}_{t,b,j}
=
\mathbf{B}_{t,b,j}^{\top}
\left(
\mathbf{d}^{\mathrm{pred}}_{t,b,j}
-
\mathbf{d}^{\mathrm{obs}}_{t,b,j}
\right).
\end{equation}
This two-dimensional residual preserves the tangential degrees of
freedom of the ray discrepancy and provides a first-order approximation
of angular error in the small-residual regime. The observed ray and its
tangent basis are updated together with the camera parameters during
optimization.

We additionally define a normalized hybrid residual that balances the
pixel reprojection residual $\mathbf{r}^{\mathrm{px}}_{t,b,j}$ and the
ray residual:
\begin{equation}
\mathbf{r}^{\mathrm{hyb}}_{t,b,j}
=
\begin{bmatrix}
\sqrt{1-\lambda}\,
\mathbf{r}^{\mathrm{px}}_{t,b,j}/s_{\mathrm{px}}
\\
\sqrt{\lambda}\,
\mathbf{r}^{\mathrm{ray}}_{t,b,j}/s_{\mathrm{ray}}
\end{bmatrix},
\qquad
\lambda\in[0,1],
\end{equation}
where $s_{\mathrm{px}}$ and $s_{\mathrm{ray}}$ are fixed normalization
scales computed from the median residual magnitudes before final
refinement. The standard pixel objective is used by default, while the
ray-domain and hybrid objectives are provided as optional final
refinement alternatives.

\section{USING THE TEMPLATE}

\begin{table*}[t]
\centering
\caption{
Cross-backend verification on the held-out multi-board and independent
checkerboard test sets. BabelCalib uses the same detected observations
and target geometry as ours.
}
\label{tab:cross_backend_verification}

\footnotesize
\setlength{\tabcolsep}{3.2pt}
\renewcommand{\arraystretch}{1.08}

\begin{tabularx}{\textwidth}{
@{}
l
>{\raggedright\arraybackslash}p{3.75cm}
@{\hspace{1pt}}
c
c
>{\centering\arraybackslash}X
c
c
@{}
}
\toprule
\multirow{2}{*}{Model}
& \multirow{2}{*}{Method / Observations}
& \multirow{2}{*}{$f_u / f_v$}
& \multirow{2}{*}{$c_u / c_v$}
& \multirow{2}{*}{Model Params.}
& \multicolumn{2}{c}{Test RMSE [px] $\downarrow$} \\
\cmidrule(lr){6-7}
&
&
&
&
& Multi-board
& Checkerboard \\
\midrule

\multirow{3}{*}{DS}
& Ours
& 1175.65 / 1175.29
& 2242.88 / 2275.61
& $[-0.1905,\ 0.6171]$
& \textbf{0.5228}
& \textbf{0.5918} \\

& BabelCalib w/ Outer+Internal
& 1176.56 / 1175.92
& 2242.90 / 2275.63
& $[-0.1914,\ 0.6175]$
& 0.5293
& 0.6020 \\

& BabelCalib w/ Outer4 only
& 1943.49 / 1943.08
& 2242.91 / 2275.67
& $[0.3401,\ 0.8048]$
& $0.5256^{\dagger}$
& -- \\

\midrule

\multirow{3}{*}{KB}
& Ours
& 1451.58 / 1451.10
& 2242.37 / 2275.38
& $[-0.0197,\ 0.0010,\ -0.0032,\ 0.0006]$
& \textbf{0.5216}
& \textbf{0.5960} \\

& BabelCalib w/ Outer+Internal
& 1454.52 / 1453.74
& 2242.91 / 2275.65
& $[-0.0196,\ -0.0018,\ -0.0007,\ -0.0001]$
& 0.5286
& 0.6052 \\

& BabelCalib w/ Outer4 only
& --
& --
& --
& --
& -- \\

\bottomrule
\end{tabularx}

\vspace{2pt}
\parbox{\textwidth}{
\scriptsize
$^\dagger$ The Outer4-only setting achieves a comparable multi-board
test RMSE but converges to a substantially shifted intrinsic parameterization.
}
\end{table*}

\begin{table}[t]
\centering
\caption{Internal-observation ablation under a fixed split and backend. Max. error denotes the maximum frame--board RMSE.}
\label{tab:internal_seed_ablation}
\footnotesize
\setlength{\tabcolsep}{1.5pt}
\renewcommand{\arraystretch}{1.05}
\begin{tabular}{@{}lcccc@{}}
\toprule
Method
& \shortstack{Reproj.\\[-1pt]Err. [px] $\downarrow$}
& \shortstack{Outer\\[-1pt]Err. [px] $\downarrow$}
& \shortstack{Internal\\[-1pt]Err. [px] $\downarrow$}
& \shortstack{Max.\\[-1pt]Err. [px] $\downarrow$} \\
\midrule
Outer-only & 0.5251 & \textbf{0.1002} & 0.5582 & 1.679 \\
Local Pinhole & 4.3647 & 0.1089 & 4.6499 & 50.353 \\
\textbf{Ours} & \textbf{0.5243} & 0.1030 & \textbf{0.5572} & \textbf{1.678} \\
\bottomrule
\end{tabular}
\end{table}

Use this sample document as your LaTeX source file to create your document. Save this file as {\bf root.tex}. You have to make sure to use the cls file that came with this distribution. If you use a different style file, you cannot expect to get required margins. Note also that when you are creating your out PDF file, the source file is only part of the equation. {\it Your \TeX\ $\rightarrow$ PDF filter determines the output file size. Even if you make all the specifications to output a letter file in the source - if your filter is set to produce A4, you will only get A4 output. }

It is impossible to account for all possible situation, one would encounter using \TeX. If you are using multiple \TeX\ files you must make sure that the ``MAIN`` source file is called root.tex - this is particularly important if your conference is using PaperPlaza's built in \TeX\ to PDF conversion tool.

\subsection{Headings, etc}

Text heads organize the topics on a relational, hierarchical basis. For example, the paper title is the primary text head because all subsequent material relates and elaborates on this one topic. If there are two or more sub-topics, the next level head (uppercase Roman numerals) should be used and, conversely, if there are not at least two sub-topics, then no subheads should be introduced. Styles named ÒHeading 1Ó, ÒHeading 2Ó, ÒHeading 3Ó, and ÒHeading 4Ó are prescribed.

\subsection{Figures and Tables}

Positioning Figures and Tables: Place figures and tables at the top and bottom of columns. Avoid placing them in the middle of columns. Large figures and tables may span across both columns. Figure captions should be below the figures; table heads should appear above the tables. Insert figures and tables after they are cited in the text. Use the abbreviation ÒFig. 1Ó, even at the beginning of a sentence.

% Table 4: single-column residual-objective evaluation.
% Generated by scripts/generate_residual_mode_multiset_table.py.
% Values are read from Stage5 result summaries; do not edit manually.
\begin{table}[t]
\centering
\caption{Residual-objective ablation under the Double Sphere model.}
\label{tab:residual_mode_multiset}
\small
\setlength{\tabcolsep}{1.6pt}
\renewcommand{\arraystretch}{1.03}
\begin{tabular}{@{}lrrrr@{}}
\toprule
\multicolumn{5}{@{}l}{\textit{(a) Holdout image-plane accuracy}} \\
Residual & Seq. A & Seq. B & Seq. C & Mean \\
 & \multicolumn{4}{c}{Reprojection RMSE [px] $\downarrow$} \\
\midrule
Pixel & 0.5193 & 0.5042 & 0.5269 & 0.5168 \\
Spherical & 0.5195 & \textbf{0.5035} & 0.5269 & 0.5167 \\
Hybrid & \textbf{0.5191} & 0.5042 & \textbf{0.5266} & \textbf{0.5166} \\
\addlinespace[1pt]
Kalibr (DS) & 0.6284 & 0.5666 & 0.6386 & 0.6112 \\
\midrule
\multicolumn{5}{@{}l}{\textit{(b) View-angle-conditioned ray accuracy}} \\
Residual & $0$--$30^{\circ}$ & $30$--$50^{\circ}$ & $\geq 50^{\circ}$ & Overall \\
 & \multicolumn{4}{c}{True angular RMSE [mrad] $\downarrow$} \\
\midrule
Pixel & 0.3204 & 0.3493 & 0.4117 & 0.3488 \\
Spherical & 0.3205 & \textbf{0.3489} & \textbf{0.4107} & \textbf{0.3486} \\
Hybrid & \textbf{0.3202} & 0.3491 & 0.4114 & 0.3486 \\
\midrule
\multicolumn{5}{@{}l}{\textit{(c) Cross-sequence geometric stability and cost}} \\
Residual & Overall & Outer & Cross-seq. & BA time \\
 & RMSE [px] & RMSE [px] & ray RMS [deg] & [s] \\
\midrule
Pixel & 0.5168 & \textbf{0.1255} & 0.2334 & \textbf{1.0} \\
Spherical & 0.5167 & 0.1275 & \textbf{0.2237} & 10.3 \\
Hybrid & \textbf{0.5166} & 0.1277 & 0.2281 & 12.6 \\
\bottomrule
\end{tabular}
\vspace{2pt}
\parbox{\columnwidth}{\footnotesize \textit{Protocol.} Polar bins are assigned once using the shared initialization camera. Cross-seq. ray RMS is the pairwise ray-curve disagreement over 1,422 shared image locations (4,266 ray pairs). BA time is the mean incremental-optimization runtime. All metrics are lower-is-better.}
\end{table}



   \begin{figure}[thpb]
      \centering
      \framebox{\parbox{3in}{We suggest that you use a text box to insert a graphic (which is ideally a 300 dpi TIFF or EPS file, with all fonts embedded) because, in an document, this method is somewhat more stable than directly inserting a picture.
}}
      %\includegraphics[scale=1.0]{figurefile}
      \caption{Inductance of oscillation winding on amorphous
       magnetic core versus DC bias magnetic field}
      \label{figurelabel}
   \end{figure}
   

Figure Labels: Use 8 point Times New Roman for Figure labels. Use words rather than symbols or abbreviations when writing Figure axis labels to avoid confusing the reader. As an example, write the quantity ÒMagnetizationÓ, or ÒMagnetization, MÓ, not just ÒMÓ. If including units in the label, present them within parentheses. Do not label axes only with units. In the example, write ÒMagnetization (A/m)Ó or ÒMagnetization {A[m(1)]}Ó, not just ÒA/mÓ. Do not label axes with a ratio of quantities and units. For example, write ÒTemperature (K)Ó, not ÒTemperature/K.Ó

\section{CONCLUSIONS}

A conclusion section is not required. Although a conclusion may review the main points of the paper, do not replicate the abstract as the conclusion. A conclusion might elaborate on the importance of the work or suggest applications and extensions. 

\addtolength{\textheight}{-12cm}   % This command serves to balance the column lengths
                                  % on the last page of the document manually. It shortens
                                  % the textheight of the last page by a suitable amount.
                                  % This command does not take effect until the next page
                                  % so it should come on the page before the last. Make
                                  % sure that you do not shorten the textheight too much.

%%%%%%%%%%%%%%%%%%%%%%%%%%%%%%%%%%%%%%%%%%%%%%%%%%%%%%%%%%%%%%%%%%%%%%%%%%%%%%%%



%%%%%%%%%%%%%%%%%%%%%%%%%%%%%%%%%%%%%%%%%%%%%%%%%%%%%%%%%%%%%%%%%%%%%%%%%%%%%%%%



%%%%%%%%%%%%%%%%%%%%%%%%%%%%%%%%%%%%%%%%%%%%%%%%%%%%%%%%%%%%%%%%%%%%%%%%%%%%%%%%
\section*{APPENDIX}

Appendixes should appear before the acknowledgment.

\section*{ACKNOWLEDGMENT}

The preferred spelling of the word ÒacknowledgmentÓ in America is without an ÒeÓ after the ÒgÓ. Avoid the stilted expression, ÒOne of us (R. B. G.) thanks . . .Ó  Instead, try ÒR. B. G. thanksÓ. Put sponsor acknowledgments in the unnumbered footnote on the first page.



%%%%%%%%%%%%%%%%%%%%%%%%%%%%%%%%%%%%%%%%%%%%%%%%%%%%%%%%%%%%%%%%%%%%%%%%%%%%%%%%

References are important to the reader; therefore, each citation must be complete and correct. If at all possible, references should be commonly available publications.



\begin{thebibliography}{99}

\bibitem{c1} G. O. Young, ÒSynthetic structure of industrial plastics (Book style with paper title and editor),Ó 	in Plastics, 2nd ed. vol. 3, J. Peters, Ed.  New York: McGraw-Hill, 1964, pp. 15Ð64.





\end{thebibliography}




\end{document}
